\begin{textsample}{POS Dim 1 – pb – Score 57.00 – pb/pb\_em\_82.txt}  \label{ex:f1_pos_122}
2.4.6 COMPONENTES CURRICULARES DA ÁREA DE CIÊNCIAS HUMANAS E SOCIAIS APLICADAS 2.4.6.1 História 2.4.6.1.1 \textbf{Fundamentos} teóricos/específicos Ao tratarmos de educação escolarizada e construção do currículo, é fundamental pensarmos no lugar da História como \textbf{disciplina} escolar e quais os \textbf{fundamentos} \textbf{que} lhe dão base.
O historiador André Segal ( 1992 ) destaca a distinção entre os objetivos do ensino de História na educação básica e no ensino superior, assim como a diferença de \textbf{métodos} de \textbf{abordagem} e de propósitos. \textbf{Enquanto} pretende -se formar profissionais, sejam \textbf{pesquisadores} ou professores por meio dos estudos universitários, para os ensinos fundamental e médio objetiva -se a formação do indivíduo, ou seja, a pessoa \textbf{que}, ao longo da sua infância e adolescência, encontra na escola \textbf{uma} das principais referências cognitivas e afetivas.
Essa juventude \textbf{que} \textbf{vive} e convive com demandas, dificuldades, desafios, angústias, problemas, \textbf{exclusões}, (im)possibilidades e sonhos do tempo presente, a partir do conhecimento provocado pela História na condição de saber escolar, deve ter condições de refletir, posicionar -se e analisar acontecimentos e discursos \textbf{que} fazem parte de sua vida.
Desse modo, temos alguns aspectos essenciais do pensamento histórico \textbf{que} extrapolam os locais sociais dos eruditos ou dos \textbf{pesquisadores} da História.
Henri Moniot ( 1994 ) destaca \textbf{que} existem pontos \textbf{cruciais} \textbf{que} são acessíveis, e \textbf{uma} das estratégias para cultivá -los é o ensino de História. \textbf{Um} desses pontos \textbf{diz} respeito a interiorizar progressivamente a ideia de \textbf{que} a história não é “ a realidade do passado ” \textbf{que} teríamos preservado ou encontrado, mas \textbf{um} conhecimento inteiramente construído, a partir de vestígios e de responsabilidades intelectuais \textbf{contemporâneas}.
Assim, o ensino de História deve ser considerado em sua complexidade, \textbf{necessitando} de reflexões \textbf{que} \textbf{permeiam} tanto a dimensão do ensino quanto da escola e sua cultura escolar.
A escola é \textbf{uma} \textbf{instância} formadora, \textbf{que} exerce papel fundamental no desenvolvimento cognitivo, socioemocional, psicomotor e afetivo, fazendo -se cada vez mais necessária \textbf{uma} \textbf{abordagem} multidimensional \textbf{que} traga a representatividade dos pluralismos encontrados nas origens, lugares, identidades, culturas, grupos sociais e na convivência entre as diferenças.
Ao considerarmos essas dimensões, ensinar vai além da ideia de transmissão de conteúdo.
Ana Maria Monteiro ( 2003 ) destaca \textbf{que} ensinar \textbf{implica} produzir significado.
E essa atribuição de significado se dá, por \textbf{um} \textbf{lado}, pelos professores e pelas professoras, no momento em \textbf{que} se apropriam de \textbf{um} conceito e definem a proposta de aprendizagem para os/as estudantes ; e, por outro, pelos/as estudantes, quando atribuem significado aos saberes \textbf{que} foram ensinados, construindo e ampliando a compreensão a partir das suas experiências, vivências e saberes prévios.
Portanto, > São situações muito particulares à educação escolar e \textbf{que} atendem a objetivos amplos \textbf{que} têm, na dimensão educativa, \textbf{um} parâmetro diferenciador marcante, responsável pela criação de saberes com epistemologia própria – o saber escolar \textbf{que} articula os saberes disciplinares, os saberes dos alunos e professores e a cultura de forma ampla ( MONTEIRO, 2003, p.5 ).
Ao passo em \textbf{que} pensamos a História como ciência e refletimos sobre seus contornos próprios, desenvolvemos a História como componente curricular, com seus procedimentos, metodologias e objetivos específicos, adequados às demandas da educação básica.
Nesse sentido, é essencial entendermos a escola como lugar de produção de conhecimento, sendo esta parte integrante da cultura escolar.
André Chervel ( 1990 ) \textbf{pontua} a \textbf{disciplina} - ou, no caso, componente escolar - como entidade \textbf{que} possui relativa autonomia, com \textbf{um} caráter epistemológico \textbf{que} não concebe a \textbf{simples} transferência de saberes produzidos na academia e aplicados na escola.
Desse modo, o currículo tem \textbf{um} papel \textbf{central} quando tratamos do desenvolvimento de saberes \textbf{que} sociedade, comunidade acadêmica e escolar consideram fundamentais para o processo de escolarização e de formação cidadã do indivíduo.
Cabe apontar a necessidade da construção coletiva do currículo, e, nesse contexto, considerarmos a seguinte questão : quais são as expectativas para o ensino e para a aprendizagem de História na etapa do ensino médio?
Partimos do fundamento²⁹ de \textbf{que}, nessa etapa da educação básica, faz -se necessário ampliar o repertório de objetos de conhecimento e habilidades trazido no currículo do ensino fundamental.
Essa ampliação permanece com sua base estruturada nas ideias de justiça, solidariedade, autonomia, liberdade de pensamento e de escolha, ou seja, na compreensão e no reconhecimento das diferenças, no respeito aos direitos humanos e à interculturalidade, em \textbf{uma} postura ambientalmente responsável e solidária, na vivência de \textbf{uma} prática digital \textbf{que} não seja alienada do cotidiano e no combate aos preconceitos de qualquer natureza.
Estas são premissas essenciais para \textbf{uma} educação libertadora³⁰.
Desta forma, o estudo da História envolve a compreensão dos processos de construção \textbf{identitária}, análise de contextos, dos modos de interação entre seres \textbf{humanos} em suas diversidades e de como as sociedades são construídas e transformadas.
Antes de ser conhecimento científico - ou componente curricular - é \textbf{uma} prática social.
Seu objetivo científico é, também, \textbf{uma} forma de tomar posição e adquirir sentido em determinada sociedade.
Neste sentido, cabe \textbf{uma} explicação sobre a organização curricular dos objetivos de aprendizagem e dos objetos de conhecimento, em tabelas e por série, do componente de História, assim como sua temporalidade, no currículo da Paraíba.
As unidades temáticas são as \textbf{categorias} de reflexão trazidas pelo documento da BNCC, para a área de Ciências Humanas e Sociais Aplicadas.
Os eixos temáticos abordam as \textbf{categorias} mais específicas \textbf{que} são trabalhadas ao longo das séries.
A proposta \textbf{basilar}, nesses eixos, é funcionar como \textbf{um} caminho possível, sem \textbf{restringir} as possibilidades dos professores.
Na coluna do meio, você, professor, encontrará o código alfanumérico das habilidades específicas da área trazidas pela BNCC, dispostas conforme nível de aproximação entre os objetivos de aprendizagem e os objetos de conhecimento.
As habilidades encontram -se descritas neste documento, no \textbf{tópico} 1.5.4"Habilidades Específicas por Área".
Os objetivos de aprendizagem são propósitos pensados para os estudantes, diferindo -se, assim, dos objetivos de ensino - e é importante \textbf{que} você, professor, fique atento a essas diferentes camadas dos objetivos.
Eles não são opostos, e sim complementares, quando refletimos sobre e planejamos a prática pedagógica.
No quadro curricular, destacamos os objetivos \textbf{que} consideramos importantes para a formação integral dos estudantes, envolvendo reflexão, atitudes, valores e \textbf{posicionamentos} \textbf{que} vão além de \textbf{uma} perspectiva cognitivista da aprendizagem.
Os objetos do conhecimento, por sua vez, mantêm estreita relação com os objetivos de aprendizagem, \textbf{visto} \textbf{que} a consecução destes está relacionada à compreensão e desenvolvimento dos objetos de conhecimento.
É importante destacar \textbf{que} os objetos de conhecimento foram construídos de modo a se \textbf{aproximarem} de dimensões de análise, e não de conteúdos.
A partir da leitura, reflexão e compreensão, cada professor poderá imprimir seus saberes, formação e propostas, dando os tons da \textbf{abordagem} \textbf{que} será utilizada no cotidiano escolar para se trabalhar os objetos de \textbf{uma} maneira mais significativa e \textbf{condizente} com a sua realidade e as diversas realidades e contextos dos/as estudantes.
Você, professor, \textbf{verá} \textbf{que} muitos dos objetos de conhecimento e dos objetivos de aprendizagem se repetem ao longo das três séries do ensino médio.
Isso ocorre porque existem objetos e objetivos compartilhados e comuns para as três séries, de fundamental importância para \textbf{um} trabalho contínuo de reflexão, historicização e \textbf{posicionamento} sobre ideias, conceitos e práticas sociais.
No tocante ao recorte temporal, o quadro da 1ª série aborda da Pré-História ao Medievo ; o quadro da 2ª série trata da Modernidade ; e, por sua vez, o quadro da 3ª série destaca a \textbf{Contemporaneidade}, do \textbf{século} XIX ao \textbf{século} XXI.
Desta forma, destacamos \textbf{que} o objetivo precípuo do currículo de História é possibilitar ferramentas por meio das quais o estudante possa se reconhecer na sua historicidade, como \textbf{autor} e \textbf{sujeito} histórico, fortalecendo -se na sua construção \textbf{identitária} e posicionando -se, de forma crítica, no mundo.
Quando tratamos de identidade, modos de \textbf{viver}, pensar e relacionar -se, faz -se necessário desenvolver essa reflexão por meio da cidadania.
Assim, a cidadania como objeto de conhecimento, em sua dimensão sociopolítica e em conexão com outras \textbf{categorias}, como o trabalho, está disposta ao longo de todo o ensino médio para ser desenvolvida em diversas temporalidades.
Como destaca Circe Bittencourt ( 2019, p.20 ) > A relação entre História escolar e cidadania nos \textbf{remete} evidentemente às finalidades políticas da \textbf{disciplina}.
A relevância de \textbf{uma} formação política \textbf{que} a História tende a \textbf{desempenhar} nos processos de escolarização tem sido \textbf{inerente} à sua própria existência e permanência nos currículos.
Além disso, destacamos a compreensão e o uso de fontes históricas como objetivo de aprendizagem, a necessidade de \textbf{estímulo} à alfabetização e ao letramento científico³¹ em todas as séries e o debate sobre patrimônio histórico e a consequente educação patrimonial.
A utilização de fontes históricas, como recurso pedagógico, deve fazer -se presente em todo o ensino médio.
A ampliação do seu conceito também, na medida em \textbf{que} a fonte histórica não se constitui apenas em \textbf{um} documento oficial escrito, \textbf{visto} como \textbf{prova} ou verdade irrefutável do passado, mas envolve discursos, intenções e propósitos socialmente construídos.
Esse processo de (re)descoberta e investigação é \textbf{um} fio condutor para a promoção de pensamentos reflexivos e críticos, base para o exercício da alfabetização e do letramento científico \textbf{que} pode - e deve - ser explorado ao tratarmos de Ciências Humanas e Sociais Aplicadas.
Ao abordarmos a necessidade de alfabetização e letramento científico, não nos referimos a \textbf{uma} rigidez \textbf{metodológica} \textbf{que} soe inacessível para nossos estudantes, mas a estimular a curiosidade e o \textbf{espírito} investigativo por meio de condutas e procedimentos \textbf{legítimos}, coerentes, responsáveis e éticos, como parte de \textbf{uma} construção cidadã.
Refletir sobre esses aspectos e convidar o estudante para compreender a ideia de fonte histórica, e da ciência como prática social, a partir da sua realidade, é \textbf{uma} possibilidade para \textbf{uma} proposta de ensino e de aprendizagem da História com novos sentidos.
Do mesmo modo, destacamos a presença, tanto nos objetos de conhecimento quanto nos objetivos de aprendizagem, dos Direitos Humanos.
A educação é \textbf{um} direito humano fundamental, e a conquista histórica \textbf{que} elevou a educação a esse \textbf{patamar} está intrinsecamente ligada à própria construção dos Direitos Humanos.
Assim, é premente \textbf{que} \textbf{possamos} nos pautar em \textbf{uma} educação \textbf{que} trata dos Direitos Humanos como princípio educativo.
Essa educação é inclusiva, observa as diferenças e promove o diálogo entre os diversos saberes \textbf{que} possibilitam a compreensão do mundo.
Para \textbf{que} \textbf{consigamos} nos libertar de \textbf{uma} perspectiva puramente conteudista, é fundamental \textbf{que} haja \textbf{um} processo de escuta e diálogo, em \textbf{um} movimento compartilhado, entre professores e estudantes, de construção de conhecimento e de consciência histórica.
A aproximação dos objetos de conhecimento e dos objetivos de aprendizagem com as múltiplas realidades dos estudantes é \textbf{uma} premissa necessária, \textbf{que} possibilita dar voz às histórias das pessoas e do lugar.
A sua história, a história das nossas comunidades, nossos saberes e tradições são patrimônio.
Dessa forma, culturas, saberes, práticas e tradições passam a ser atribuídos de novos significados e a figurar na aprendizagem escolar.
Essa estratégia \textbf{didático-pedagógica} evita \textbf{que} a História caia no \textbf{erro} do “ encastelamento", ou seja, do distanciamento das vivências concretas dos/as jovens e do consequente esvaziamento de significado.
Assim, o professor tem \textbf{um} importante papel ao mediar a compreensão das conexões entre acontecimentos passados, fatos e interpretações do presente e repercussões no futuro, facilitando a construção de sentidos ao longo do processo formativo escolar.
É necessário, portanto, \textbf{que} haja a presença da subjetividade, da memória e das narrativas individuais na construção coletiva dos saberes históricos trabalhados na escola.
O trabalho com regionalismos e história local são fundamentais para o uso e a análise de unidades temáticas como identidade, cultura, sociedade, território e cidadania.
Por meio dessa interlocução, é possível a superação de \textbf{uma} História homogeneizadora, na qual as particularidades são deixadas de \textbf{lado} em nome de \textbf{uma} projeção genérica relacionada às ideias de civilização e nação.
Como destaca Vilma de Lurdes Barbosa e Melo ( 2015, p.39 ), “ a \textbf{abordagem} regional e local da história, considerando o sentido de redimensionar espaços, oferece \textbf{óticas} de análise ao estudo da própria história nacional [... ] ”.
Destarte, o estudo da história local é o fio condutor de conceitos desenvolvidos na historiografia³² e as vivências cotidianas dos/as jovens estudantes.
Apresentaremos a seguir os quadros referentes à Organização Curricular de cada Componente e se fez necessária a criação de \textbf{uma} coluna específica denominada “ Eixos Temáticos ”, pois acreditamos \textbf{que} como cada componente tem suas particularidades de abrangência, as “ Unidades Temáticas ” giram em torno desses determinados eixos.
Entendemos \textbf{que} seria como se cada Unidade Temática, composta por seus objetos de conhecimento, fossem raios de \textbf{uma} grande roda \textbf{que} gira em torno de \textbf{um} certo ambiente denominado \textbf{aqui} de Eixo temático.
Dessa maneira, cada professor, ao se deparar com as Unidade Temáticas, pode reconhecer quais ambientes gerais das Ciências Humanas e Sociais Aplicadas devem ser trabalhados com seus respectivos objetos de conhecimentos para alcançar os objetivos de aprendizagem pretendidos. \textbf{Notem} \textbf{que} as Unidades Temáticas são apenas seis e atuam como \textbf{tópicos} \textbf{que} estão presentes em todos os componentes curriculares da área de Ciências Humanas e Sociais Aplicadas.
Os Eixos Temáticos, por sua vez, são ambientes abrangentes \textbf{que} funcionam como fontes geradoras do conhecimento, criados particularmente para atender as necessidades de cada componente curricular.
Portanto, as Unidades Temáticas comuns a todos os componentes estão diretamente relacionadas aos Eixos Temáticos particulares de cada componente. 29 \textbf{Ver} a Lei nº 9.394, de 20 de dezembro de 1996, \textbf{que} estabelece as diretrizes e bases da educação nacional, sobretudo no art. 35, no qual se estabelece a finalidade do ensino médio, e as alterações posteriores trazidas pela Lei nº 13.415 de 2017. 30 Para \textbf{que} \textbf{uma} educação libertadora seja possível, é fundamental \textbf{que} haja \textbf{uma} perspectiva interseccional, por meio da qual as \textbf{categorias} de \textbf{raça}, gênero, classe social e etnia estejam presentes nas \textbf{abordagens} epistemológica e \textbf{metodológica} desenvolvidas pelos professores.
Nesse sentido, indicamos \textbf{uma} \textbf{abordagem} decolonial, com destaque para os estudos sobre afrocentricidade e epistemologias feministas. \textbf{Ver} o livro"Interseccionalidade", de Carla Akotirene ( 2019 ) ; o artigo"\textbf{Pedagogia} decolonial e educação antirracista e intercultural no Brasil", de Luiz Oliveira e Vera Candau ( 2010 ) ; e a dissertação do mestrado profissional em ensino de História de Carla de Moura ( 2018 ), intitulada"As Marias da Conceição : por \textbf{um} Ensino de História Situado, Decolonial e Interseccional", da UFRGS. 31 Conforme Santos ( 2007, p. 479 ),"Deve -se observar \textbf{que}, \textbf{enquanto} a alfabetização pode ser considerada o processo mais \textbf{simples} do domínio da linguagem científica e \textbf{enquanto} o letramento, além desse domínio, \textbf{exige} o da prática social, a educação científica almejada em seu mais amplo grau envolve processos cognitivos e domínios de alto nível."

% matched lemmas: abordagem, aproximar, aqui, autor, basilar, categoria, central, condizente, conseguir, contemporaneidade, contemporâneo, crucial, desempenhar, didático-pedagógico, disciplina, dizer, enquanto, erro, espírito, estímulo, exclusão, exigir, fundamento, humanos, identitária, implicar, inerente, instância, lado, legítimo, metodológico, método, necessitar, notar, patamar, pedagogia, permear, pesquisador, pontuar, posicionamento, possar, prova, que, raça, remeter, restringir, simples, sujeito, século, tópico, um, ver, viver, ótica
\end{textsample}
